\chapter{Introduction}

Les automates cellulaires~(AC)\footnote{%
  Pour des collections th\'ematiques voir
  \cite{Demongeot85,Wolfram86a,Legendi87,Gutowitz90a}, %
  pour des ouvrages th\'ematiques voir
  \cite{Preston84,Toffoli87,Jackson91c}, %
  pour une approche didactique voir %
  \cite{Casti89a,Schroeder90a}.} %
sont des syst\`emes dynamiques virtuels, %
discrets en espace, en temps et en \'etats, homog\`enes, synchrones et
d\'eterministes. %
L'espace est constitu\'e de cellules identiques dot\'ees d'un r\'eseau
d'interconnexion r\'egulier~(une grille de dimension~2 ou~3 par
exemple). %
Une cellule est une variable discr\`ete~(une cha\^ine de bits dont la
taille d\'etermine le nombre d'\'etats de la cellule). %
Le temps est une s\'equence d'espaces g\'en\'er\'ee par l'application
synchrone d'une fonction de transition d\'eterministe. %
La fonction de transition est locale dans le temps et dans l'espace,
son argument est un ensemble fini d'\'etats de cellules voisines du
r\'eseau d'interconnexion, pour un nombre fini de points ant\'erieurs. %

Les AC sont sujets d'\'etudes ou d'applications dans les domaines des
math\'ematiques %
\cite[]{Jen90,Vichniac90}, %
de l'architecture des ordinateurs %
\cite[]{Legendi87a,Zsoter87}, %
de la programmation parall\`ele %
\cite[]{Preston84,Almasi87} %
et de la simulation de syst\`emes naturels %
tant physiques\footnote{%
  En particulier les \aquote{Lattice Gas Automata}~\cite[]{Despain90,Margolus90}.} %
\cite[]{Margolus84,Toffoli84,Rossi93} %
que biologiques\footnote{%
  L'\'etude de la vie artificielle inspire de nombreux travaux sur les
  AC, souvent assez \'eloign\'es d'observations biologiques. %
  Une exception notable \'etant \cite{Hameroff89}, qui propose
  l'existence d'une capacit\'e au calcul cellulaire pour le
  cytosquelette.} %
\cite[]{Tamayo89,Boer93}. %
Ils constituent un champ d'\'etudes des math\'ematiques discr\`etes et un
mod\`ele important en th\'eorie du calcul\footnote{%
  Voir par exemple \cite{Lindgren90} pour une \'equivalence
  entre certains AC~1D et une machine de Turing.}, %
en particulier pour les probl\`emes li\'es au parall\'elisme\footnote{%
  Les AC math\'ematiques mettent en jeu
  un espace infini, en particulier \`a limites p\'eriodiques. %
  Notre point de vue est \'egalement empirique et
  impl\'ementatoire: nous viserons sp\'ecifiquement la r\'ealisation de
  programmes dot\'es du meilleur rendement. %
  En cons\'equence, dans un contexte de ressources m\'emoire limit\'ees, nos
  AC seront \'evidemment dot\'es d'un espace restreint, par d\'efaut torique.}; %
leur structure est tr\`es proche de l'architecture des machines
SIMD\footnote{%
  \cite{Greussay85} indique que la structure des AC est tr\`es proche de
  l'architecture des machines SIMD, et nous pensons que le progr\`es des
  connaissances et des outils dans ces deux domaines sont intimement
  li\'es. %
  Les travaux de \cite{Rozin87,Rozin88,Rozin93} sur la programmation
  plane et de \cite{Signorini92,signorini94} sur la programmation par
  configurations cellulaires t\'emoignent de la fertilit\'e de cette
  approche. %
  La puissance de calcul et la taille de la m\'emoire disponible sur des
  machines aujourd'hui largement accessibles permet d'envisager la
  r\'ealisation de multiples exp\'eriences dans le domaine du parall\'elisme
  synchrone massif, sans attendre l'arriv\'ee de machines SIMD\@. %
  Les AC fournissent un outil conceptuel en mesure de jouer un r\^ole
  central dans la recherche d'objets naturels~(comportement \'emergent)
  ou artificiels~(comportement construit), pouvant servir de base \`a la
  d\'efinition d'outils pratiques d'exploitation des ressources du
  parall\'elisme synchrone massif.}; %
ind\'ependamment de l'architecture des machines, de nombreux
algorithmes se pr\^etent \`a une formulation
cellulaire~\cite[]{Uhr87,Uhr84}, en particulier les algorithmes de
traitement d'images~\cite[]{Preston84,Hasselbring92}; enfin, les AC
peuvent \^etre vus comme une version discr\`ete de syst\`emes
d'\'equations diff\'erentielles\footnote{%
  Et \`a ce titre permettent la mod\'elisation de nombreux syst\`emes
  naturels. %
  La mod\'elisation d'un syst\`eme naturel par un ensemble d'\'equations
  diff\'erentielles \'etablit des relations diff\'erentielles entre les
  variables macroscopiques du syst\`eme~(des populations par exemple). %
  La mod\'elisation par AC \'etablit des relations entre les composants
  discrets \'el\'ementaires~(les cellules) du syst\`eme. %
  Les variables macroscopiques deviennent des propri\'et\'es \'emergentes du
  syst\`eme dont le choix peut \^etre fait a posteriori, mais le mod\`ele
  doit imp\'erativement reposer sur une description en termes de
  comportement collectif du ph\'enom\`ene \'etudi\'e. %

  La description d'un syst\`eme naturel en termes de comportement
  collectif n'implique pas l'identification d'une cellule \`a un
  individu. %
  Par exemple, dans une simulation de dynamique des populations, une
  cellule peut \^etre identifi\'ee \`a une densit\'e de population en un point
  d'espace. %
  N\'eanmoins la validit\'e de la mod\'elisation~(conformit\'e de la dynamique
  des populations d'un AC \`a un syst\`eme d'\'equations diff\'erentielles
  standard de r\'ef\'erence par exemple) d\'epend toujours de la pr\'esence
  d'un nombre minimum de cellules dans la simulation.}. %

\section[Objectifs et contraintes]{Nos objectifs et contraintes pour un
  environnement de d\'e\-ve\-lop\-pe\-ment et d'observation d'AC}

Nous cherchons \`a d\'efinir et \`a r\'ealiser un environnement de
d\'eveloppement et d'observation de ces syst\`emes. %
Une des motivations principales de ce travail r\'eside dans le d\'esir de
voir~(et montrer) les s\'equences anim\'ees g\'en\'er\'ees par un AC\@. %
En effet s'il est ais\'e de construire des AC en vue de donner des
instantan\'es photographiques d'\'etats successifs, il est plus difficile de
les visualiser en continu filmique. %
Bien qu'il existe des outils\footnote{%
  \mysoftcite[Cellsim]{cellsim}, %
  \mysoftcite{scamper}, %
  \mysoftcite[CA~LAB]{ca_lab}~\cite[]{Rucker89}, %
  \mysoftcite{xlife}, %
  \mysoftcite[HODGE-C]{hodge_c}, %
  \mysoftcite{cellang}. %
  Les num\'eros en indice sup\'erieur entre crochets renvoient au r\'epertoire des
  outils, donn\'e en appendice.} %
de sp\'ecification d'AC et d'observation de leur dynamique, aucun de
ceux dont nous connaissons l'existence ne r\'epond
pleinement aux contraintes que nous allons d\'efinir. %
Nous souhaitons \'egalement que cet outil permette imm\'ediatement la
recherche et la cr\'eation de nouveaux AC, puis la constitution d'un
catalogue d'AC connus. %
Enfin nous esp\'erons qu'il serve \`a l'identification d'AC pouvant servir
de primitives dans un futur langage bas\'e sur la combinaison d'AC. %
Les qualit\'es recherch\'ees sont la rapidit\'e, l'extensibilit\'e, ainsi que la
facilit\'e de sp\'ecification et d'interaction\footnote{%
  La \aconcept{rapidit\'e} peut \^etre d\'efinie comme le rendement d'une
  impl\'ementation   pour une architecture donn\'ee. %
  Dans ce contexte, et pour une architecture cible s\'equentielle, la
  rapidit\'e s'exprime comme le nombre d'instructions
  \'el\'ementaires de la machine cible utilis\'ee pour simuler une
  instruction de la machine source. %
  Cette d\'efinition est n\'eanmoins trop g\'en\'erale, en effet, si nous
  proposons dans le cadre de ce travail des moteurs cellulaires \`a
  rendement constant~(\ie ind\'ependant de l'espace initial et de la
  r\`egle de transition), il est possible de concevoir des moteurs
  cellulaires \`a rendement variable: \mysoftcite{xlife} d\'ecoupe
  l'espace cellulaire en blocs et permet le calcul sur des
  configurations cellulaires virtuelles de grandes tailles;
  \cite{Gosper84} d\'ecrit un algorithme bas\'e sur une m\'emoire cache
  asynchrone des configurations cellulaires~(le cache contient un
  treillis de morceaux d'espace-temps cellulaire). %

  L'extensibilit\'e peut, dans notre contexte, se ramener \`a la
  conception d'un ensemble d'objets constituant \`a la fois la bo\^ite \`a
  outils et les \'el\'ements de base de la construction des exp\'eriences
  cellulaires. %
  La recherche de cette qualit\'e peut entrer en conflit avec la
  recherche de la rapidit\'e, par exemple les variations de param\`etres
  sur les moteurs cellulaires \`a rendement constant que nous avons
  r\'ealis\'es passent par une phase de macro-g\'en\'eration de code source.}. %


\subsection*{Le cadre syst\`eme, mat\'eriel, et logiciel:~les performances}

Notre environnement de d\'eveloppement est Unix et X11, notre
architecture cible est une machine standard de~10~\`a~100~mips
et~8~\`a~64 \astrangeterm{m\'egabytes}{%
  Nous emploierons indiff\'eremment les termes octet et byte, ainsi que
  leurs d\'eriv\'es: kilobyte, m\'egaoctet ou gigaoctet.} %
de m\'emoire utile. %
Pour un AC bidimensionnel compos\'e de \asquare{256}
cellules de~2~bits\footnote{%
  Avec un voisinage de Moore~(9~bits de voisinage), 2~bits d'\'etats
  donnent une table de taille 512~kilobit~($2^{2 \times  9} \times  2$), ce qui
  est tr\`es loin des limites que nous nous sommes fix\'es, mais le passage
  \`a trois bits de voisinage donne une table de taille
  48~m\'egaoctets~($\frac{2^{3 \times  9} \times  3}{8}$) qui se trouve \`a la limite
  des ressources m\'emoires consid\'er\'ees. %
  Nous insistons cependant sur le fait que rien n'interdit en pratique
  l'usage de telles tables, m\^eme si la connaissance de la
  distribution des configurations locales devient n\'ecessaire pour
  pr\'edire les performances du syst\`eme constitu\'e du triplet
  \emph{table} de transition, \emph{configuration} initiale,
  \emph{performance} intrins\`eque du syst\`eme de gestion m\'emoire de la
  machine supportant l'application.} %
avec un voisinage de Moore, le syst\`eme doit \^etre capable de calculer
une dizaine d'\'etats par seconde~(10~Hz)\footnote{%
  D'un point de vue \aquote{esth\'etique}, cette fr\'equence correspond \`a
  un minimum subjectif permettant de restituer l'illusion du mouvement
  continu~(les images t\'el\'evisuelles ont une fr\'equence de~25~Hz, les
  images cin\'ematographiques modernes une fr\'equence de~24~Hz, les
  images cin\'ematographiques originelles une fr\'equence de~14 \`a~20~Hz,
  le passage \`a des fr\'equences plus \'elev\'ees est dict\'e par des
  consid\'erations techniques li\'ees aux besoins du parlant pour le cin\'ema
  et \`a la fr\'equence du secteur alternatif pour la t\'el\'evision). %
  Cette fr\'equence est donc \'elev\'ee, car d'un point de vue
  \aquote{syst\`eme} elle correspond~(pour notre espace de r\'ef\'erence
  $256^2$) \`a un d\'ebit de~640~kilo \bialt{op\'erations}{cellules} par
  seconde, ce qui pour une machine hypoth\'etique \`a~20~mips, correspond
  \`a~32~instructions de la machine cible pour une instruction de la
  machine source~(une mise \`a jour cellulaire).}. %
Le co\^ut li\'e \`a la visualisation est plus difficile \`a \'evaluer, mais la
cha\^ine de traitement compl\`ete~(production, extraction, visualisation)
doit permettre l'observation d'un film \`a une fr\'equence sup\'erieure
\`a~1~Hz\footnote{%
  Eisenstein, Gance ou Godard par leur utilisation occasionnelle de
  l'alternance de plans \`a une telle fr\'equence illustrent le proc\'ed\'e
  esth\'etique. %

  Cette fr\'equence correspond par ailleurs \`a un imp\'eratif ergonomique
  subjectif concernant le \aquote{temps de r\'eponse} d'une application,
  elle borne la taille des espaces accessibles \`a un mode de
  visualisation s\'equentielle lorsque la contrainte de l'illusion du
  mouvement est lev\'ee. %
  
  La perte de l'illusion du mouvement, n'est pas contradictoire \`a
  l'\'etude du mouvement. %
  Le physiologiste Marey \'etudie le vol des oiseaux:~son %
  \aquote{fusil photographique}~(1882) et son %
  \aquote{Chronophotographe \`a pellicule}~(1888) %
  servent \`a l'analyse photographique du mouvement. %
  Bien que consid\'er\'e comme pr\'ecurseur du cin\'ematographe,
  il ne se pr\'eoccupe pas de restituer l'illusion du mouvement. %
  En cin\'ema d'animation, la \anewterm{pixillation} utilise sciemment
  la s\'equence contre l'illusion. %
  \begin{quote}
    ``L'animation n'est pas l'art des dessins qui bougent, mais l'art des
    mouvements dessin\'es. %
    Ce qui ce passe entre chaque image a beaucoup plus d'importance que
    ce que l'on voit sur chaque image.'' \textsc{N. McLaren}.{\hfill$\cdots/\cdots$}%
  \end{quote}
  En 1941 dans ``L'\'evolution cr\'eatrice'', Bergson
  d\'efend l'id\'ee que ``le m\'ecanisme de notre connaissance usuelle est
  de nature cin\'ematographique''~\cite[page 305]{Bergson89}. %
  Pour lui ``la forme n'est qu'un instantan\'e pris sur une
  transition'', seul est r\'eel le changement continuel de forme. %
  Il place cette m\'etaphore au c{\oe}ur d'une r\'eflexion historique qui
  remonte \`a Z\'enon d'\'El\'ee. %

  Dans son acharnement \`a d\'efendre ``L'ant\'eriorit\'e ontologique du
  continu sur le discret'', \cite{Thom92} va jusqu'\`a proposer
  une ``version informatis\'ee de la preuve cart\'esienne de
  l'existence de Dieu'' d\'eclarant contradictoire la finitude du nombre
  de neurones c\'er\'ebraux et la pens\'ee de l'infini. %
  Sans pr\'etendre trancher une question aussi d\'ebattue, nous pouvons
  n\'eanmoins sugg\'erer que l'\'etude des AC propose une approche \`a rebours
  de cette question, la recherche pratique de l'immanence
  ph\'enom\'enologique du continu rempla\c{c}ant les tentatives de
  d\'emonstration de son ant\'eriorit\'e ontologique. %

  Sur notre configuration de d\'eveloppement~(un~486DX2 avec~32~m\'egaoctets
  de m\'emoire, le syst\`eme d'exploitation Linux et X11 pour la
  visualisation) la fr\'equence typique pour une cha\^ine comprenant la
  g\'en\'eration sur un voisinage de Moore~(toujours pour un espace de
  r\'ef\'erence $256^2$) et l'extraction d'un plan m\'emoire~(constitu\'e d'un
  des deux bits d'\'etat) dans un premier processus, pip\'e sur un processus
  dumpant les plans cellulaires vers le serveur X11~(tournant sur la
  m\^eme machine) est sup\'erieure \`a~10~Hz. %
  Le d\'ecouplage des diff\'erentes op\'erations de la cha\^ine permet
  d'obtenir une fr\'equence de g\'en\'eration sup\'erieure \`a~15~Hz~(module de
  g\'en\'eration redirig\'e sur \texttt{/dev/null}) et une fr\'equence de
  visualisation monochrome~(pour un seul plan de bits cellulaire)
  sup\'erieure \`a~30~Hz, (mesur\'e via un fichier historique
  de~8~m\'egaoctets pour~1024~g\'en\'erations). %
  Ces mesures s'entendent en temps \'ecoul\'e r\'eel, pour un syst\`eme
  autrement au repos, les mesures de performance de visualisation ne
  constituent \'evidemment pas une r\'ef\'erence \`a notre travail, mais elles
  permettent de consid\'erer, au regard des mesures de performance du
  moteur cellulaire, que celui-ci atteint le d\'ebit plafond d'un moteur
  \`a rendement constant, \ie la meilleure performance possible est
  toujours garantie dans le pire des
  cas~(espace cellulaire et portion d'espace-temps satur\'e).}. %
Nous avons ainsi r\'ealis\'e un noyau d'objets permettant la construction
d'exp\'eriences, sous la forme d'une librairie C, interfac\'ee \`a
\asoft{Tcl}\footnote{%
  Tool Command Language.}%
~\cite[]{Ousterhout90,Ousterhout91}. %

\subsection*{Sp\'ecification:~la description des exp\'eriences cellulaires}

La sp\'ecification concerne les r\'eseaux d'interconnexions, les fonc\-tions
de transitions, et le s\'equencement des op\'erations applicables sur
l'espace. %
Les fonc\-tions de transitions sont \'ecrites en~\asoft{C}, et sont dot\'ees
d'une convention d'interface pr\'evue pour permettre une ind\'ependance
des param\`etres \`a leur localisation dans la m\'emoire d'\'etats d'une
cellule et la r\'eutilisation dans des fonctions de plus haut niveau. %
Les r\'eseaux d'interconnexions sont d\'efinis et cr\'e\'es en~\asoft{Tcl} \`a
partir des primitives de la librairie. %
Les espaces sont cr\'e\'es en~\asoft{Tcl}. %
Le s\'equencement est d\'efini par un programme~\asoft{Tcl}. %

\subsection*{Interaction:~l'interface humaine}

L'interaction se fait soit par l'interm\'ediaire d'un shell~\asoft{Tcl}, soit
par une interface X11 bas\'e sur~\asoft{Tk}. %

\subsection*{Visualisation:~les outils graphiques}

La visualisation se fait par l'interm\'ediaire de fen\^etres X11
pr\'esentant des s\'equences d'espaces monochromes~(plans de bits de
l'espace cellulaire) ou couleur~(tranches de bits de l'espace
cellulaire). %
Les variables macroscopiques~(\'el\'ements statistiques) sont visualis\'ees
par l'interm\'ediaire d'objets X11 sp\'ecialis\'es~(plotters). %

\mydumpfull[htb]{casNew}%
{Une interface humaine exp\'erimentale pour le Workbench}%
{Une interface humaine exp\'erimentale et son dumper}%
{Une interface X11 r\'ealis\'ee avec
  \protect\mysoftcite{Tk}~\cite[]{Ousterhout90,Ousterhout91}. %
  Le dumper pr\'esente une image typique de la transformation d'un espace
  initial al\'eatoire de densit\'e\,\Sundemi par la
  r\`egle~\arule{Anneal}~\cite[]{Vichniac86,Toffoli87}\footnote{%
    \mygetrule[phrase]{anneal}} %
  Le dumper pr\'esente \'egalement diff\'erents compteurs li\'es aux temps
  propre de la r\`egle, du s\'equenceur, et au pool d'espaces de l'animation. %
}

\section{Les outils existants:~pas de standard}

Une recherche aussi exhaustive que possible des outils publics
existants nous a convaincu de la n\'ecessit\'e d'une utilisation critique
de certains de ces outils, en particulier dans le domaine de la
visualisation, associ\'e \`a une r\'ealisation exp\'erimentale ad hoc
d'\'el\'ements d'architecture de calcul cellulaire. %

Les outils de d\'eveloppement et de visualisation d'automates
cellulaires sont peu nombreux\footnote{%
  Bien que les r\'ealisations de \arule{Life} soit
  innombrables, et les packages cellulaires plus ou moins g\'en\'eraux tr\`es
  nombreux, leurs r\'ealisations minimales posant peu de probl\`emes, nous ne
  prenons en compte que les packages dont il existe des annonces
  explicites.}. %
Le plus connu d'entre eux est sans doute \mysoftcite{cellsim};
\mysoftcite{scamper} et \mysoftcite[CA~LAB]{ca_lab} sont \'egalement des
outils disponibles depuis 1989. %
Il existe aussi des outils sp\'ecialis\'es pour un mod\`ele d'automate
cellulaire. %
\mysoftcite{xlife} fournit une r\'ealisation d'ex\'ecution tr\`es rapide du
jeu de la vie~\cite[]{Berlekamp82}, \mysoftcite[HODGE-C]{hodge_c} une
r\'ealisation de la machine hodge-podge de
Gerhard~\&~Schuster~\cite[]{Dewdney88}, une classe d'AC inspir\'ee par
les r\'eactions chimiques autocatalytiques, telle que la r\'eaction de
Beluso\mbox{v-Z}habotinski. %
Enfin, \mysoftcite{cellang}, d'inspiration plus architecturale,
propose un langage de description d'AC~\cite[]{Eckart92a}. %
\nocite{Eckart92b}

\section[Le calcul et la visualisation]%
{Un nouvel environnement:~exploration par le calcul et la visualisation}

Nous avons donc r\'ealis\'e un nouvel environnement\footnote{%
  Ci-apr\`es d\'esign\'e par le terme \anewterm{workbench}.} %
de d\'eveloppement et de visualisation d'automates cellulaires. %

Nous destinons ce workbench tant \`a l'exploration de la dynamique
des automates cellulaires~(AC), qu'\`a la construction et \`a l'\'evaluation
d'architectures et d'algorithmes SIMD\@. %
De plus nous proposons de consid\'erer la r\'ealisation d'algorithmes SIMD
sur machines SISD non pas comme un palliatif \`a l'absence de machines
SIMD, mais comme un mode d'utilisation particulier des ressources
disponibles sur les machines SISD modernes~(50~Mips, 32~m\'egaoctets). %

La programmation des calculateurs massivement parall\`eles r\'eels et
virtuels pose le double probl\`eme du pouvoir d'expression et de
l'efficacit\'e d'un langage au regard d'une architecture. %
Ce probl\`eme se pose \'egalement pour les architectures SISD, mais l'on
peut consid\'erer qu'il concerne d\'esormais l'ad\'equation des langages aux
applications plut\^ot qu'aux architectures. %
Optant pour une approche r\'esolument bottom-up de la recherche des
primitives, nous proposons d'utiliser les AC comme composants
\'el\'ementaires d'une combinatoire visant \`a la construction d'objets
\bialt{cellulaires}{SIMD} arbitrairement complexes. %
Ces objets se divisent en deux classes \'epist\'emologiques. %
Nous les nommerons objets \aconcept{analytiques} et objets
\aconcept{synth\'etiques}. %
Les objets analytiques sont typiquement issus de l'application
it\'erative d'une fonction simple sur un espace cellulaire initial
al\'eatoire. %
Ils recoupent intuitivement le champ d'\'etude traditionnel des AC\@. %
Les objets synth\'etiques sont issus de la s\'election d'une fonction
localement universelle, et de son application contr\^ol\'ee~(globalement
param\'etr\'ee) sur un espace initial soigneusement configur\'e. %
Ils correspondent aux architectures des machines SIMD\@. %
De l'\'etude des objets analytiques nous esp\'erons tirer des
enseignements permettant \`a terme de consid\'erer les processus
morphog\'en\'etiques comme des \'el\'ements de programmation. %
De la construction d'un outillage SIMD nous attendons les moyens de
contr\^ole n\'ecessaires \`a l'\'etude des objets analytiques. %
Cet objectif, certes ambitieux, repose pr\'ecis\'ement sur la r\'ealisation
par \aconcept{lookup-tables} des fonctions de bas niveau manipulables
et combinables au sein de l'environnement. %

Notre \'etude est organis\'ee autour de la pr\'esentation des structures de
donn\'ees choisies pour le workbench; ces structures de donn\'ees
repr\'esentent \aconcept{l'espace}, \aconcept{le temps}, et
\aconcept{les lois}:~objets \'el\'ementaires du workbench. %
Nous proposons des exemples d'op\'erations applicables sur ces objets. %
Nous proposons \'egalement de consid\'erer la r\'ealisation de fonctions par
lookup-tables comme une structure de donn\'ees fondamentale, tant pour
des raisons d'efficacit\'e, que de g\'en\'eralit\'e. %

Nous consid\'erons la visualisation comme un service primordial, qui
propose l'usage du workbench en tant que laboratoire de manipulation
d'espace cellulaire~(l'espace, mais aussi les fonctions et les
dynamiques). %
Nous pr\'esentons plusieurs exemples de visualisations principalement
associ\'ees \`a des objets analytiques. %
Nous montrons comment peut s'effectuer le passage d'une m\'ethode
naturaliste \`a une m\'ethode artefactuelle en proposant une r\'ealisation
par lookup-table d'une architecture SIMD sp\'ecifique, et une m\'ethode de
transformation d'une lookup-table quelconque en s\'equence de contr\^ole
pour cette machine. %

Constatant l'absence d'outils adapt\'es aux contraintes de la
visualisation des espace-temps cellulaires\footnote{%
  La visualisation de volume de donn\'ees dans un format inadapt\'e
  conduit \`a la saturation inutile des capacit\'es m\'emoire de la machine
  d'accueil d'un calcul cellulaire.}, %
constatant \'egalement l'existence d'espace-temps cellulaires binaires a
priori toriques, nous avons con\c{c}u, r\'ealis\'e et test\'e une biblioth\`eque
de manipulation d'espace\footnote{%
  Coupe, s\'equences de coupes sur des axes de projection variables, %
  r\'eduction int\'egrale sur un axe arbitraire pond\'er\'e par la distance \`a
  la surface de projection pour obtenir des effets de transparence.} %
cellulaires aux fins de visualisation. %
Sur cette base, nous avons par exemple pu tester l'id\'ee d'une
visualisation relativiste des espace-temps:~%
nous avons projet\'e le c\^one de la vitesse de la lumi\`ere\footnote{%
  La vitesse de la lumi\`ere d'un espace-temps cellulaire est le rayon
  du voisinage de la famille de r\`egle \`a laquelle appartient la r\`egle
  g\'en\'eratrice de l'espace-temps.} %
d'un espace-temps~2D sur un espace~2D, r\'esumant ainsi sur un seul
espace~2D la dynamique de la r\`egle visualis\'ee.

\mydumpfull[p]{superglider}%
{Un glider superluminique}%
{Un glider subluminique et un glider superluminique}%
{Nous pr\'esentons dans cette figure la d\'ecouverte
  d'un \aquote{glider superluminique}. %
  Cette s\'erie d'espaces de configurations cellulaires est une s\'equence
  de tranches d'espace-temps non align\'ees sur l'axe du temps. %
  Malgr\'e la distorsion introduite par ce d\'ecalage, la r\`egle
  g\'en\'eratrice de cette portion d'espace-temps est
  parfaitement identifiable:~\arule{Life}\footnote{%
    Les fl\`eches en haut \`a droite des deux premi\`eres photos pointent
    une configuration cyclique de~\arule{life} reconnaissables.}. %
  A premi\`ere vue, la cons\'equence de l'utilisation d'un axe non
  standard pour la projection de la dynamique de~\arule{Life} est
  d'ordre purement spatial\footnote{%
    Cette appr\'eciation subjective n'est applicable qu'\`a l'exp\'erience
    de la vision du \aquote{film}.}, %
  exception faite du d\'ecalage introduit par le d\'esaxage temporel. %
  Le carr\'e sur les photos pointe vers un objet \arule{Life} de type
  glider, dot\'e d'une vitesse de d\'eplacement subluminique. %
  L'observation de l'espace-temps contenant la portion pr\'esent\'ee au
  moyen de films \`a haute fr\'equence\footnote{%
    Sup\'erieure \`a la fr\'equence t\'el\'evisuelle.} %
  nous a permis d'identifier \`a la vue un objet paradoxal, point\'e par
  un cercle sur les photos:~un glider superluminique.}

\noindent La figure~\mygetdump{superglider} pr\'esente un des usages
possible de notre librairie de manipulation d'espaces dans le
contexte du workbench:~%
nous y pr\'esentons la d\'ecouverte d'un \aconcept{glider
  superluminique}\footnote{%
  nous emploierons le terme glider~(planeur) de mani\`ere g\'en\'erique, \ie
  pour d\'esigner des objets cellulaires p\'eriodiques, \`a un d\'ecalage spatial pr\`es.}, %
en d\'eplacement \`a une vitesse apparente double de la
vitesse de la lumi\`ere. %
Nous attendions \'evidemment de cette exp\'erience des aberrations
temporelles, des illusions ou mirages inex\-plicables dans le contexte
d'une projection parall\`ele \`a l'axe du temps. %
N\'eanmoins les premiers r\'esultats faisaient plut\^ot appara\^itre une
relative invariance de la perception visuelle subjective associ\'ee \`a la
r\`egle. %
La proc\'edure utilis\'ee produit un d\'eplacement apparent des
configurations, mais les objets se comportent par ailleurs, c'est \`a
dire relativement \`a l'espace apparent pour l'exp\'erimentateur, de
mani\`ere subluminique ou luminique usuelle. %
Ce glider, le seul que nous ayons trouv\'e~(dans l'espace-temps
englobant la portion pr\'esent\'ee) \'etait donc pr\'evisible, puisqu'il
correspond \`a un glider en mouvement corr\'el\'e au front de coupe, mais
n\'eanmoins totalement surprenant. %
Beaucoup moins pr\'evisibles nous apparaissent les cons\'equences
g\'en\'erales d'observations exploitant syst\'ematiquement des projections
d'espace-temps non parall\`eles \`a l'axe du temps.

\myfigfullR[htb]{.75}{timepyrcub}%
{Tranches d'espace-temps et pyramides relativistes}%
{Sch\'emas de r\'ealisation de tranches d'espace-temps et de pyramides relativistes}%
{Le cube, sur ces sch\'emas, repr\'esente une portion d'espace-temps pour
  un AC~2D. %
  L'axe du temps est indiqu\'e par la fl\`eche verticale, %
  les pointill\'es larges sont des ar\^etes de plans de coupe. %
  Le sch\'ema de gauche illustre la proc\'edure utilis\'ee pour r\'ealiser les
  s\'equences de coupes pr\'esent\'ees en figure~%
  \mygetdump{superglider}, %
  \mygetdump{tsum-128x3-diag}, %
  et \mygetdump{anneal-256x3-diag}. %
  Les coupes sont calcul\'ees en \aquote{distance de Manhattan}, %
  Nous avons con\c{c}u un calcul d'intersection bas\'e sur une
  version~3D de l'algorithme de trac\'e de ligne de Bresenham\footnote{%
    Nous avons modifi\'e pour cet usage la version de Bob
    Pendleton~(1992) d\'eriv\'ee de~\cite{Heckbert90}.}. %
  Le sch\'ema de droite illustre la proc\'edure utilis\'ee pour r\'ealiser les
  projections relativistes\footnote{%
    La projection est construite par \aconcept{cercles de voisinage}
    remontant le temps, dont le rayon d\'etermine l'\^age, \'eventuellement
    associ\'es \`a une translation spatiale figur\'ee par la fl\`eche horizontale
    pleine.} %
  pr\'esent\'ees en figures~%
  \mygetdump{life-256x3-rel}, %
  \mygetdump{life-256x3-relmov}, %
  \mygetdump{tsum-128x3-relmovedge}, %
  }

\mydumpfull[p]{tsum-128x3-diag}%
{Une s\'equence de tranches d'espace-temps}%
{Une s\'equence de tranches d'espace-temps de la r\`egle~\arule{Qmean}}%
{La r\`egle~\arule{Qmean} est pr\'esent\'ee, ainsi que les autres r\`egles
  utilis\'ees pour la synth\`ese des images r\'ealis\'ees
  en~\mygetref[name]{the-rules}. %
  La figure~\mygetfig{timepyrcub} illustre la proc\'edure de r\'ealisation
  de cette s\'equence. %
  D'autres visualisations associ\'ees \`a la r\`egle~\arule{Qmean} sont
  pr\'esent\'ees par les figures~%
  \mygetdump{tsum-grey}, %
  \mygetdump{tsum}, %
  \mygetdump{tsum-ray-one} %
  et \mygetdump{tsum-ray-all}.}

\mydumpfull[p]{life-256x3-rel}%
{Projection relativiste fixe}%
{Une projection relativiste pour une cellule fixe}%
{Une projection relativiste~\mygetfig[paren]{timepyrcub} d'un
  espace-temps de \arule{life} pour une cellule de r\'ef\'erence fixe,
  r\'ealis\'ee sur un espace-temps~\acube{256}, pour une densit\'e
  initiale\,\Sundemi. %
  La cellule de r\'ef\'erence est au centre de la projection, %
  celle-ci appara\^it donc au sein d'un espace \aquote{\^ag\'e},
  compos\'e en plus grand nombre d'objets point-fixe et %
  de densit\'e moyenne  plus faible %
  que l'espace situ\'e \`a proximit\'e des bords de la projection, %
  de densit\'e moyenne plus \'elev\'ee et peupl\'e d'objets plus jeunes\footnote{%
    Que l'on peut cependant qualifier d'archa\"iques du point de vue de
    l'origine de la projection.}. %
  Il est possible d'accentuer l'\^age des objets lointains en diminuant la
  vitesse de la lumi\`ere utilis\'ee lors de la projection. %
  Un tel d\'ecouplage du rayon de voisinage accro\^it le probl\`eme
  d'interpr\'etation des projections produites.}

\mydumptwothirds[htb]{life-256x3-relmov}%
{Projection relativiste mobile}%
{Une projection relativiste \`a cellule mo\-bile}%
{Cette projection diff\`ere de la figure~\mygetdump{life-256x3-rel} par
  l'ajout d'une translation appliqu\'ee \`a la cellule de r\'ef\'erence.}

\mydumpfull[t]{anneal-256x3-rel}%
{}%
{Projection relativiste sur les 3~axes d'un cube d'espace-temps}%
{Cette figure regroupe trois projections relativistes de la
  r\`egle~\arule{Anneal} r\'ealis\'ees sur une m\^eme trajectoire, %
  pour l'axe temporel et les deux axes nord-sud et est-ouest. %
  L'origine du temps est visible dans le haut de la projection pour
  ces deux axes.}

\myfigfullR[htb]{.75}{transcub}%
{Transparence globale et par plan de bits}%
{Sch\'ema de r\'ealisation de r\'eductions par transparence}%
{Comme pour la figure~\mygetfig{timepyrcub}, les cubes repr\'esentent une
  portion d'espace-temps pour un AC 2D. %
  La fl\`eche pointill\'ee illustre une travers\'ee orthogonale non parall\`ele
  \`a l'axe du temps, lors d'une telle travers\'ee on peut soit r\'ealiser un
  film\footnote{%
    A d\'efaut de films, incompatible avec le support papier, nous
    pr\'esentons des s\'equences associ\'ees \`a une travers\'ee non
    orthogonale~\mygetfig[paren]{timepyrcub}.}, %
  soit une r\'eduction int\'egrale~(une somme de la valeur de toutes
  les cellules rencontr\'ees). %
  Pour les r\'eductions int\'egrales destin\'ees \`a la production d'images,
  la valeur de la cellule est pond\'er\'ee par la distance parcourue. %
  La figure~\mygetdump{tsum-128x3-trans-byte} pr\'esente une telle
  image. %
  Pour les transparences non orthogonales, les effets de bord de la
  m\'ethode utilis\'ee~(distance de Manhattan) posent un probl\`eme
  d'interpr\'etation des d\'eformations. %
  Les quatre cubes \`a droite symbolisent l'op\'eration de s\'eparations d'un
  espace-temps de cellules \`a quatre bits en quatre espace-temps de
  cellules \`a un bit. %
  La figure~\mygetdump{tsum-128x3-trans-bits} pr\'esente des
  transparences sur trois axes pour cinq \aquote{plans} de bits
  de~\arule{Qmean}, %
  La figure~\mygetdump{life-256x3-trans-bits} pour l'unique plan d'un
  espace-temps de~\arule{life}.}

\mydumpfull[htb]{anneal-256x3-diag}%
{}%
{Une s\'equence de tranches d'espace-temps de la r\`egle~\arule{Anneal}}%
{Les premi\`eres coupes correspondent \`a un temps proche du temps initial,
  mais chaque tranche pr\'esente des caract\'eristiques archa\"iques plus ou
  moins marqu\'ees.}

\mydumpfull[htb]{tsum-128x3-faces}%
{}%
{Faces du cube d'espace-temps de~\arule{Qmean}}%
{Cette figure pr\'esente les six faces d'un cube d'espace-temps de la
  r\`egle~\arule{Qmean}. %
  Les deux faces situ\'ees dans l'axe du temps, ainsi que leur \^ages
  sont clairement identifiables.}

\mydumpfull[p]{tsum-128x3-trans-bits}%
{}%
{Une transparence \'eclat\'ee par plan de bits}%
{Les cinq bits de poids fort de la r\`egle~\arule{Qmean} sont utilis\'es
  respectivement pour r\'ealiser une
  transparence~\mygetfig[paren]{transcub} sur les trois axes de
  l'espace-temps. %
  Ces s\'equences sont \`a comparer \`a la
  figure~\mygetdump{tsum-128x3-trans-byte} utilisant la totalit\'e des
  bits de chaque cellule pour la transparence.}

\mydumpfull[p]{tsum-128x3-trans-byte}%
{}%
{Une transparence utilisant tout l'espace de valeur d'une cellule}%
{Cette transparence fournit une image plus habituelle d'une
  dynamique cellulaire consid\'er\'ee comme un voxmap\footnote{%
    Le voxel est la version 3D du pixel.} %
  Des exp\'eriences ant\'erieures men\'ees avec des packages de visualisation
  scientifique \aquote{classiques}\footnote{%
    Au sens ou ils travaillent n\'ecessairement en flottants, ou consid\`erent
    des donn\'ees dont la localisation n'est pas implicite
    et en tout cas ne sont pas adapt\'es \`a un travail au niveau du bit. %
    Nous avons recherch\'e, port\'e et exp\'eriment\'e les outils suivants:~%
    \mysoftcite[XDataSlice]{XDS} et son format %
    \mysoftcite{HDF}, %
    \mysoftcite{VIS-5D}~\cite[]{Hibbard94}, %
    \mysoftcite{viewit} et %
    \mysoftcite{LIC}~\cite[]{Cabral93}.} %
  donnent pour ce genre de visualisation des r\'esultats proches en terme
  d'images produites\footnote{%
    Mais pas en termes de ressources consomm\'ees pour la production. %
    Un autre inconv\'enient de ces packages, \'etant les erreurs
    introduites par la conversion flottante interm\'ediaire lorsque
    l'espace 2D de transparence produit est destin\'e \`a la r\'ealisation
    de mesures statistiques, autant que d'images.} %
  de ceux pr\'esent\'es ici, %
  mais l'usage de nos fonc\-tions de manipulation d'espaces cellulaires
  est n\'ecessaire\footnote{%
    D'un point de vue pratique, la conversion d'un bitmap
    tridimensionnel en voxmap de flottants, \'eventuellement reconverti
    en voxmap de bytes, est en effet parfaitement inad\'equate.} %
  \`a la pro\-duc\-tion des transparences par plan de bits de
  la figure~\mygetdump{tsum-128x3-trans-bits}. %
  La comparaison des images produites par les m\'ethodes plan de bits
  \vs voxel, sugg\`ere l'\'etude des relations entre les deux
  formulations, et une possible g\'en\'eralisation de la premi\`ere \`a
  l'analyse de voxmaps.} %

\mydumpfull[htb]{anneal-256x3-trans-bits}%
{}%
{Une transparence par plan de bits pour~\arule{Anneal}}%
{La m\'ethode de production comment\'ee en
  figure~\mygetfig{transcub} est utilis\'ee 
  pour la r\`egle~\arule{Anneal}. %
  Le nombre d'\'etats de~\arule{Anneal} \'etant \'egal \`a deux, la m\'ethode
  \'eclat\'ee se confond avec la m\'ethode standard.}

\mydumpfull[htb]{life-256x3-trans-bits}%
{}%
{Une transparence par plan de bits pour~\arule{Life}}%
{La m\'ethode de production comment\'ee en
  figure~\mygetfig{transcub} est utilis\'ee 
  pour la r\`egle~\arule{Life}}

\mydumpfull[htb]{tsum-128x3-relmovedge}%
{}%
{Une d\'etection de contour appliqu\'ee sur une projection relativiste}%
{L'accumulation de transformations renforce le probl\`eme de
  l'interpr\'etation des images produites. %
  Ces figures illustrent les limites rencontr\'ees en combinant une
  d\'etection de contour une \`a projection relativiste \`a cellule de
  r\'ef\'erence mobile, appliqu\'ee sur les trois axes de projection d'un
  cube d'espace-temps de la r\`egle~{Qmean}.}

L'\'etat actuel du d\'eveloppement du workbench n'a pas encore permis
l'exp\'erimentation de constructions complexes d'objets mixtes. %
La mesure de variables d'\'etats macroscopiques, par exemple, utilise
actuellement des outils externes, au lieu d'algorithmes SIMD adapt\'es \`a
la nature discr\`ete et finie des espaces cellulaires. %

%%
%% Local Variables:
%% mode: auto-fill
%% iso-accents-minor-mode: nil
%% TeX-command-default: "mytex"
%% TeX-master: "top"
%% Local IspellParsing: "latex-mode ~latin1"
%% Local IspellPersDict: "~/.ispell_lisp"
%% LocalWords: "d\'evelop pement d'AC cellsim scamper CA LAB ca lab xlife hodge"
%% LocalWords: "HODGE-C cellang impl\'ementation d'espace-temps mips m\'egabytes"
%% LocalWords: "byte kilobyte m\'egaoctet gigaoctet bidimensionnel Moore kilobit"
%% LocalWords: "Eisenstein Gance Godard Marey Chronophotographe pixillation ur"
%% LocalWords: "McLaren Bergson Z\'enon d'\'El\'ee finitude ph\'enom\'enologique DX Linux"
%% LocalWords: "pip\'e dumpant dev null monochrome interfac\'ee Tcl Tool Command htb"
%% LocalWords: "Language s\'equencement casNew Workbench dumper Tk Anneal shell ad"
%% LocalWords: "s\'equenceur monochromes plotters hoc packages hodge-podge Gerhard"
%% LocalWords: "Schuster Beluso v-Z macro-g\'en\'eration habotinski workbench SIMD"
%% LocalWords: "SISD bottom-up lookup-tables manipulables artefactuelle NEW tsum"
%% LocalWords: "superluminique subluminique luminique timepyrcub relativiste rel"
%% LocalWords: "diag anneal Manhattan Bresenham Pendleton relmov relmovedge XDS"
%% LocalWords: "Qmean the-rules tsum-grey tsum-ray-one tsum-ray-all paren voxmap"
%% LocalWords: "est-ouest transcub trans-byte trans-bits voxel XDataSlice HDF"
%% LocalWords: "viewit LIC bitmap bytes versus voxmaps nord-sud name"
%% LocalWords: "lookup-table relativistes glider point-fixe"
%% LocalWords: "espace-temps l'espace-temps superglider d\'esaxage"
%% End:

%% Objectifs et contraintes
%% Outils existants
%% Calcul et visualisation

\section{Pourquoi notre nouvel environnement?}

Tout d'abord, m\^eme si l'histoire de l'\'etude des AC est presque aussi
ancienne que l'histoire des ordinateurs~\cite[]{Metropolis93b}, les
architectures cellulaires construites, bien que non encore
industrialis\'ees apparaissent comme les plus prometteuses\footnote{%
  Il est difficile de savoir si ce fait est \`a mettre au cr\'edit des
  architectures cellulaires, nous nous bornerons \`a proposer l'id\'ee que
  la perception du potentiel d'usage d'une architecture est
  inversement proportionnelle \`a sa r\'ealisation industrielle.}, %
et le probl\`eme des m\'ethodes de constructions d'architectures et
d'algorithmes cellulaires restent largement ouvert. %
De m\^eme, \`a propos du probl\`eme de
l'auto-reproduction~\cite[]{vonNeumann66}, la dualit\'e opposant
analytique \`a synth\'etique accompagne la naissance des AC\footnote{%
  Voir \cite{Reggia93} pour une approche permettant d'accorder plus
  d'autonomie aux ph\'enom\`enes d'auto-reproduction. %
  Nous consid\'erons, dans ce contexte, l'autonomie comme la mesure
  d'une position sur l'axe qui va du construit \`a l'\'emergeant.}. %
Cette dualit\'e, ou plut\^ot la capacit\'e offerte par les AC de basculer
d'un univers \'epist\'emologique \`a un autre tout en restant au sein d'un
mod\`ele unique, doit \^etre incarn\'ee dans l'environnement. %
Nous souhaitons donc un environnement permettant une approche mixte. %

Par ailleurs, aucun des outils que nous avons cit\'es ne r\'epond
totalement \`a un ensemble de contraintes plus pratiques que nous avons
fix\'ees comme cadre de d\'eveloppement. %
\begin{itemize}
\item la nature publique\footnote{%
    Par \aconcept{publique} nous entendons accessible sous forme
    source et utilisable sans licence commerciale.} %
  des outils
\item la rapidit\'e d'ex\'ecution
\item la richesse de visualisation
\item la souplesse de sp\'ecification
\item un environnement d'exploitation Unix et X11
\item la modularit\'e des constituants en vue de la construction et du
  s\'equencement d'exp\'eriences cellulaires. %
\end{itemize}

\asoft{CA~LAB} n'est pas un outil public et n'est pas disponible dans
l'environnement Unix. %
\asoft{Cellsim} de par sa d\'ependance \`a~\asoft{SunView} n'est
utilisable que dans un environnement Sun\footnote{%
  Felicity George a r\'ealis\'e en 1992 un portage de~\asoft{cellsim} pour
  l'environnement~\asoft{XView}.}, %
\asoft{scamper} pr\'esente une application int\'egrant un formalisme de
sp\'ecification des r\`egles bas\'ees par pattern, des compteurs de
populations et une interface visuelle bas\'ee sur X11, mais il ne permet ni
ex\'ecution rapide, ni s\'equencement. %
\asoft{Cellang} propose un langage de description d'AC, mais son
compilateur ne produit pas un code permettant de tirer le meilleur
parti des ressources disponibles. %

\section{Des AC pour l'optimisation d'usage de ressources}

Concernant le parall\'elisme massif, les tendances de la pr\'ec\'edente
d\'ecennie en architecture des machines ne confirment pas les
pr\'edictions avanc\'ees \`a son propos\footnote{%
  Machines SIMD massives reposant sur des techniques de fabrication de
  type wafer scale int\'egration, associ\'ees \`a des algorithmes de
  reconfigurations de r\'eseaux d'interconnexion permettant la
  r\'esistance aux d\'efauts; pyramides de telles architectures pour la
  vision et plus g\'en\'eralement l'IA, ou architectures diverses pour la
  physique, la biologie mol\'eculaire, le pattern matching haut d\'ebit,
  la r\'ealisation de film holographique temps r\'eel.}. %
N\'eanmoins, hors consid\'erations architecturales, les ressources
accessibles continuent \`a cro\^itre \`a un rythme soutenu\footnote{%
  Au point qu'on pourrait mettre en doute l'application du
  principe des gaz parfaits \`a la consommation de ressources. %
  Contre cette tendance, nous proposons de consid\'erer qu'il n'est plus
  temps d'apprendre comment se contenter de ressources restreintes,
  mais comment tirer profit de ressources disponibles.}. %
Hormis le calcul intensif classique, les gains de ressources semblent
s'orienter vers le parall\'elisme MIMD massif pour la gestion et le
support multim\'edia pour le poste de travail. %
Nous affirmons que les AC constituent un mod\`ele efficace et universel
pour l'utilisation de ces ressources, y compris par la prise en compte
du parall\'elisme multi-calculateur\footnote{%
  C'est \`a dire l'utilisation d'un ensemble de machines
  autonomes~(stations de travail ou autres, localement parall\`eles ou
  non) reli\'ees par un r\'eseau d'interconnexion~(de type bus ou point
  \`a point) dot\'ees d'une couche logicielle syst\`eme ad hoc, comme une
  ressource parall\`ele en soi.} %
associ\'e \`a l'interconnexion par \astrangeterm{WAN}{%
  Wide Area Network, il s'agit ici de remarquer que les d\'ebits
  disponibles ou annonc\'es pour les r\'eseaux de t\'el\'ecommunication publics
  permettent d'envisager la construction virtuelle de
  machines cellulaires reparties sur une large zone g\'eographique.} %
de ceux-ci. %
La r\'egularit\'e structurelle des donn\'ees et des traitements des AC
associ\'es \`a l'absence de probl\`eme algorithmique de synchronisation
des \'el\'ements du calcul r\'eparti, permet un calcul et une affectation de
ressource simple. %
Une approche bottom-up de la sp\'ecification du comportement des
\'el\'ements cellulaires permet, de plus, de s'affranchir des probl\`emes
d'architecture CPU, une fonction consid\'er\'ee sous l'angle de sa table
de transition est ind\'ependante des architectures de
traitement\footnote{%
  A une permutation de l'ordre des bits sur le mot pr\`es, et toujours en
  consid\'erant que l'unit\'e de traitement r\'eellement disponible reste
  fondamentalement une architecture de von Neumann dot\'ee d'une m\'emoire
  adressable importante.}. %

\mysetref{repartition-du-calcul}
A propos de la r\'epartition d'un calcul cellulaire sur r\'eseau de
stations de travail\footnote{%
  Algorithme distribu\'e avec recouvrement: calcul des bords, envoi des
  bords, calcul du centre, r\'eception des bords.}, %
et concernant le probl\`eme de l'overhead associ\'e au passage de messages
impl\'ementant la synchronisation des bords lorsque le calcul courant
r\'eclame un acc\`es \`a une cellule voisine, ind\'ependamment d'une mesure
particuli\`ere des d\'ebits respectifs du couple
\bialt{cpu}{m\'emoire}~($D_{m}$) et du canal de communication($D_{c}$),
le rapport des tailles \bialt{locales}{distantes} $N^{d}/N^{d-1} = N$
garantit l'existence d'une valeur $k$ de $N$ telle que l'usage de
plusieurs n{\oe}uds de calculs sera plus efficace que l'usage d'un
seul d'entre eux:~$\frac{D_{m}}{D_{c}} k > 1$. %
%% \[\frac{D_{m}k^{d}}{D_{c}k^{d-1}} > 1\]

A d\'efaut de machines cellulaires industrielles, nous
esp\'erons pouvoir utiliser les DSP et les acc\'el\'erateurs
\aquote{bitblt}~\cite[]{Pike84} aujourd'hui largement disponibles
comme acc\'el\'erateurs de nos moteurs cellulaires. %

\section[AC et morphog\'en\`ese]{Des AC comme outils de r\'eflexion morphog\'en\'etique}

L'essence de la forme pose un probl\`eme, ne serait ce que
par l'existence de la g\'en\'ericit\'e du moule. %
Tout \'etat du monde n'est pas une forme, toutefois une forme est par
d\'efinition une configuration d'espace cellulaire prise parmi un
ensemble fini de configurations. %
Cette d\'efinition r\'eclame un crit\`ere de
\bialt{discrimination}{r\'eduction} de l'ensemble des \'etats. %
Une premi\`ere possibilit\'e r\'eside dans l'utilisation d'une fonction
arbitraire sur un \'etat d\'eterminant la valeur d'un attribut formel:
esth\'etique, g\'eom\'etrique, sta\-tistique, arithm\'etique. %
Les formes sont figuratives, abstraites, s\'eduisantes, rondes, carr\'ees,
sym\'etriques, rares, uniques\footnote{%
  Chaque configuration d'espace est \'evidemment unique, on pourrait
  n\'eanmoins mesurer, intuitivement ou formellement, l'unicit\'e \`a partir
  du produit de l'ensemble des \'etats par l'ensemble des
  fonctions de qualification formelle disponibles. %
  Les fonctions de qualification sont des fonctions de r\'eduction de
  l'information~(le nombre de bits rendu est inf\'erieur au nombre de
  bits fourni), le cardinal de l'image d'un ensemble de formes est
  donc inf\'erieur \`a celui de l'ensemble de formes, la distribution des
  redondances de valeurs induites par la diff\'erence de cardinalit\'e
  d\'eterminant les classes d'\'equivalence de cet attribut. %
  
  Une relation d'ordre pour l'unicit\'e est alors d\'efinie par une
  fonction minimisant le cardinal de l'union des classes d'\'equivalence
  des images des formes.}, %
r\'eguli\`eres, fractales, denses, color\'ees, leur num\'ero d'ordre enfin,
peut \^etre argument d'une fonction choisie au hasard. %
Cette m\'ethode pose le probl\`eme de l'\'enum\'eration des fonctions de
r\'eduction\footnote{%
  Ainsi que celui du choix \`a priori des \'etats consid\'er\'es. %
  Les fonctions de qualification sont implicitement destin\'ees \`a
  introduire une r\'eduction du nombre des formes par l'introduction de
  classes d'\'equivalence sur l'ensemble des \'etats.}. %
Une deuxi\`eme m\'ethode consiste \`a opposer la g\'en\'eticit\'e \`a la
g\'en\'ericit\'e, c'est \`a dire l'histoire de la forme au moule de la forme. %
Dans ce contexte la r\'eduction du nombre d'espaces appel\'es \`a la forme
est d\'etermin\'ee par l'existence d'une trajectoire dans l'espace des
\'etats. %
Cette trajectoire \'etant elle m\^eme d\'etermin\'ee par l'application it\'er\'ee
d'une fonction de l'ensemble des \'etats vers l'ensemble des \'etats. %
Les trajectoires ne sont pas moins nombreuses que les \'etats, mais si
l'on ajoute une contrainte de localit\'e aux fonctions g\'en\'eratrices de
trajectoires, une tentative d'\'enum\'eration devient accessible. %
Cette contrainte de localit\'e constitue le fondement d'une approche
cellulaire de la forme. %
Cette contrainte introduit un lien entre la vitesse et la forme. %
Plus g\'en\'eralement, la notion de contrainte s'applique ais\'ement au
monde des formes dans une perspective esth\'etique, le libre choix des
contraintes de production de l'objet d\'eterminant la qualit\'e artistique
de l'{\oe}uvre, dont la forme peut \^etre con\c{c}ue comme une trace
d'ex\'ecution\footnote{%
  La perception de l'{\oe}uvre n'est jamais exempte de la perception
  de l'histoire de sa fabrication~\cite[]{Gombrich71}.}. %
De m\^eme la \aquote{forme} d'un espace cellulaire
devient la trace d'un programme\footnote{%
  De l'application d'une loi non ergodique sur un espace initial
  n\'ecessairement al\'eatoire.}. %

\section[AC, espace, temps et calcul]%
{Des AC comme outils de r\'eflexion sur l'espace, le temps et le calcul}

Nous avons d\'ej\`a mentionn\'e l'existence d'une dualit\'e
analytique-synth\'etique. %
D'autres projections de cette dualit\'e sont possibles, sur les axes
dynamique statique, heuristique algorithmique, parall\`ele s\'equentiel ou
temporel spatial. %
La trajectoire d'un AC sur un espace de configuration est un objet
dynamique, la structure d'un espace de r\`egles de transition d'AC pour
un voisinage donn\'e est un objet statique. %
A la fronti\`ere se trouvent des objets donn\'es par l'application
it\'erative d'une fonction simple sur un espace cellulaire initial
globalement param\'etr\'e~(l'automate auto-reproducteur de von Neuman, ou
une machine de Turing impl\'ement\'ee par un espace de configuration
d'un AC tel \arule{Life} de Conway). %

La quantit\'e de parall\'elisme d'une ex\'ecution peut se mesurer au nombre
d'instructions diff\'erentes rencontr\'ees lors de cette ex\'ecution. %
Il existe une \'equivalence entre un AC d\'efini par un r\'eseau
d'interconnexion et une fonction de transition, et un algorithme
cellulaire d\'efinit par une s\'equence d'instructions
cellulaires. %
Si un AC est consid\'er\'e comme une instruction cellulaire \'el\'ementaire,
celle-ci poss\`ede un co\^ut, qui s'exprime comme la taille de la table
de v\'erit\'e de cette instruction cellulaire. %
Si un AC est consid\'er\'e comme une s\'equence d'instructions cellulaires,
celle-ci poss\`ede un co\^ut, qui s'exprime comme la somme des temps de
r\'ealisation de chaque instruction, et in fine la somme de la taille
des tables de v\'erit\'e de chaque instruction \'el\'ementaire. %
L'impl\'ementation optimale d'un algorithme cellulaire serait obtenu
lorsque, pour une occupation maximale de la m\'emoire disponible pour les
tables de v\'erit\'es~(les fonctions de transition) des instructions
cellulaires, le nombre de primitives diff\'erentes composant l'algorithme
est minimum. %
Le ratio du nombre de primitives employ\'ees sur le nombre de primitives
existantes d\'etermine l'importance des effets relevant de la
dynamique~(l'heuristique) de ceux relevant de l'algorithme proprement
dit dans la fonction associ\'ee au programme cellulaire.

Nous disons que la fonction de transition d'un AC est localement
univer\-selle\footnote{%
  Nous dirons parfois qu'un AC est globalement universel, pour
  d\'esigner la propri\'et\'e usuelle d'universalit\'e par opposition \`a la
  propri\'et\'e d'universalit\'e locale des PE de machines SIMD\@. %
  L'universalit\'e locale implique \'evidement l'universalit\'e globale. %
  \arule{Life} constitue dans le contexte de ce travail l'arch\'etype
  d'un AC dot\'e de l'universalit\'e globale, et le processeur GAPP
  l'arch\'etype d'un AC dot\'e de l'universalit\'e locale.} %
si une seule cellule consid\'er\'ee comme PE se comporte comme une machine
universelle sous r\'eserve d'une extension virtuelle du nombre d'\'etats
cellulaires. %
C'est a priori le cas de toutes les architectures SIMD,
m\^eme si le faible nombre d'\'etats accessibles\footnote{%
  Ce faible nombre d'\'etats doit n\'eanmoins \^etre exprim\'e en logarithme,
  c'est \`a dire en nombre de bits adresse.} %
destine \'evidemment ces machines \`a un usage cellulaire. %

D'un cot\'e le \aquote{calcul s\'equentiel} pur d\'eroule une s\'equence
d'instructions d\'efinie en extension sur un espace initial \'egalement
d\'efini en extension, de l'autre le \aquote{calcul cellulaire pur}
it\`ere une loi unique d\'efinie en extension sur un espace initial dont
seules des propri\'et\'es statistiques\footnote{%
  La taille d'un espace cellulaire incarn\'e, n\'ecessairement fini, doit
  n\'eanmoins \^etre consid\'er\'ee comme tendant vers l'infini; la
  sp\'ecification d'une configuration cellulaire particuli\`ere en
  extension devient alors une barri\`ere \`a la pleine utilisation des
  ressources disponibles.} %
sont d\'efinies. %
La distinction entre une s\'equence d'instructions et une \emph{loi}
apparaissant comme une instruction unique, ne tient ni \`a la taille de
la loi ni \`a une propri\'et\'e telle l'universalit\'e de l'architecture
sous-tendant la s\'equence d'instructions, ou l'universalit\'e
locale ou globale de la loi\footnote{%
  En effet toute s\'equence d'instruction, tout algorithme peuvent
  \'evidemment \^etre transform\'es en une instruction
  unique par le calcul d'une lookup-table.} %
mais \`a l'opposition entre l'unicit\'e de l'application d'une s\'equence
d'instructions, d'un algorithme et la n\'ecessaire it\'eration sans
fin\footnote{%
  Comme pour l'espace, ressource n\'ecessairement limit\'ee devant
  n\'eanmoins \^etre consid\'er\'ee comme tendant vers l'infini, le temps
  cellulaire, le nombre de cycles d'it\'erations de l'application de la
  loi sur l'espace doit \^etre consid\'er\'e comme tendant vers l'infini.} %
de la loi sur l'espace cellulaire. %
A la fronti\`ere\footnote{%
  La structure de l'espace englobant le s\'equentiel et le cellulaire
  appara\^it ici tr\`es clairement.} %
le calcul cellulaire it\'erant une loi unique sur un espace dont la
configuration ini\-tiale est une transformation~(un mapping) d'une s\'erie
d'instructions d'un calcul s\'equentiel, c'est \`a dire le calcul
cellulaire exploitant l'universalit\'e globale d'une loi; ou encore le
calcul cellulaire SIMD d\'eroulant une s\'equence d'instructions d\'efinie
en extension, c'est \`a dire le calcul cellulaire utilisant
l'universalit\'e locale d'une loi, et sans doute, mais sans qu'il soit
possible d'en \'etablir simplement la part, l'universalit\'e globale de
cette loi. %

Si l'on consid\`ere que l'encodage d'un algorithme sur une machine
s\'equentielle pose un probl\`eme aigu, et encore davantage
l'encodage du probl\`eme \`a traiter, il en va bien s\^ur de m\^eme pour les
machines SIMD. %
On continue toutefois dans le contexte SIMD \`a r\'esoudre des probl\`emes
algorithmiques, \`a utiliser l'universalit\'e pour la construction de
machines algorithmes. %
Lorsque l'on passe \`a une machine cellulaire pure, la nature
intrins\`equement statistique du processus de calcul\footnote{%
  Nous avons jusque l\`a \'evoqu\'e la seule incertitude statistique
  associ\'ee \`a la sp\'ecification de la configuration initiale, la loi
  \'etant consid\'er\'ee comme purement d\'eterministe, on voit bien cependant
  comment un ind\'eterminisme sur l'espace est de fait indiff\'erentiable
  d'un ind\'eterminisme sur la loi, laquelle n'est jamais que de
  l'espace local g\'en\'erateur du temps.} %
change radicalement le point de vue. %
Le probl\`eme \`a traiter ne pouvant d\`es lors gu\`ere s'exprimer que par
analogie \`a une morphog\'en\`ese
cellulaire\footnote{%
  Expression \`a prendre dans un sens non restrictif \`a la biologie bien
  s\^ur, \`a l'existence pr\`es d'un paysage \'epig\'en\'etique.}. %
Il sera toujours possible de plonger une approche SIMD, voire
purement s\'equentielle et distribu\'ee dans un contexte cellulaire pur,
mais outre que cela pose des probl\`emes qui nous apparaissent assez
ardus de faisabilit\'e pratique\footnote{%
  Cela pose aussi des probl\`emes plus fondamentaux, dont le moindre
  n'est pas d'\'eclaircir la nature ontologique de l'incertitude
  statistique qui semble bien devoir \^etre li\'ee \`a une d\'efinition d'un
  calcul cellulaire r\'eellement massif, m\^eme si la dimension
  technologique de cette incertitude existe \'egalement.}, %
cela ne tend pas \`a r\'epondre \`a la
question de la sp\'ecificit\'e du calcul cellulaire pur. %
En effet, si celui-ci ne semble pas, dans son approche la plus stricte,
isomorphe \`a une machine universelle arbitraire, ce pourrait \^etre d\^u \`a
la nature particuli\`ere des morphog\'en\`eses cellulaires, lesquelles ne
sont peut-\^etre pas r\'eductibles \`a des algorithmes. %
Le calcul s\'equentiel peut \^etre consid\'er\'e comme une version
adimensionnelle~(de dimension z\'ero) d'un calcul cellulaire. %
Le calcul cellulaire algorithmique\footnote{%
  C'est \`a dire la r\'ealisation d'un algorithme par une une
  configuration cellulaire initiale associ\'ee \`a une loi globalement
  universelle.} %
correspondant au choix d'un point unique dans l'espace des \'etats
cellulaires, permettant par le plongement d'une machine de Turing dans
l'espace cellulaire, de r\'ealiser une r\'eduction
dimensionnelle de l'espace cellulaire. %
Le calcul cellulaire SIMD\footnote{%
  C'est \`a dire la r\'ealisation d'un algorithme sur une configuration
  cellulaire par une m\'ethode s\'equentielle locale.} %
correspondant au choix d'un mapping probl\`eme vers configuration cellulaire,
associ\'e ou non \`a l'utilisation des propri\'et\'es dynamiques d'une loi ou
plusieurs lois pour contr\^oler une trajectoire dans l'espace des
configurations cellulaires. %
Le calcul cellulaire pur, par la prise en compte de l'incertitude
statistique, cherche \`a dresser un paysage g\'en\'eral du calcul,
lequel devient une propri\'et\'e intrins\`eque, voire \'emergeante d'un mod\`ele
d'univers compos\'e d'atomes d'espace dot\'e d'une loi universelle. %
Une des propri\'et\'es particuli\`erement int\'eressante, de notre point de
vue, du calcul cellulaire, r\'eside dans les \'equivalences espace-temps
rendues possibles par son architecture. %
En effet la mesure des quantit\'es \'el\'ementaires de cette architecture,
ainsi que les transformations mutuelles possibles entre ces quantit\'es
sont d'impl\'ementation ais\'ee. %
Ainsi, en calcul cellulaire pur, le nombre d'\'etats d'une cellule et le
nombre de voisins de cette cellule d\'etermine une borne sup\'erieure \`a la
complexit\'e\footnote{%
  Nous n'avons pas de d\'efinition pr\'ecise de cette complexit\'e,
  il s'agit a priori d'une mesure statistique de la richesse de
  comportement, rapport\'ee \`a une dur\'ee. %
  La capacit\'e et la rapidit\'e morphog\'en\'etique associ\'ees \`a une r\`egle.} %
du calcul effectu\'e\footnote{%
  L'analyse de la table de transition permet de pr\'eciser une valeur
  de complexit\'e, en particulier la distance \`a la borne sup\'erieure. %
  Une table constante, les tables identit\'e ou inverse sur un seul
  bit du mot de voisinage~(translation, et translation alternance de
  la configuration initiale) sont faciles \`a identifier. %
  Une g\'en\'eralisation possible de ces notions conduit \`a la d\'efinition
  de param\`etres de contr\^ole de l'espace des r\`egles tels que le
  param\`etre $\lambda$ de Langton ou le param\`etre d'ench\^assement.}. %
La taille de \arule{Life} est de~512~bits,
celle de l'automate auto-reproducteur de von Neumann de
$2^{5\lceil\log_{2}(29)\rceil}\lceil\log_{2}(29)\rceil =
20 \text{m\'egaoctets}$. %
Il serait int\'eressant de conna\^itre la taille minimale d'une
configuration initiale de \arule{Life} r\'ealisant une machine
universelle, il n'existe pas \`a notre connaissance de r\'ealisation
effective d'une telle configuration. %
La taille de la configuration r\'ealisant l'automate auto-reproducteur
de von Neumann a \'et\'e \'evalu\'ee \`a plusieurs centaines de milliers de
cellules. %
%% John kemeny 1955
Jaqueline Signorini a \'etudi\'e les conditions de r\'ealisation de cet
automate sur le GAPP~\cite[]{Signorini92}, mais une configuration
auto-reproductrice compl\`ete n'a encore jamais \'et\'e men\'ee \`a bien. %

L'impl\'ementation de \arule{Life} sur une machine s\'equentielle
$M_{\text{SISD}}$ utilise $a$ cycles de base pour chacune des $N$
cellules, et $b = 2$ bits\footnote{% 
  Un plan \avar{past} et un plan \avar{futur}, impl\'ementant le
  synchronisme de mise \`a jour, on peut envisager des optimisations
  r\'eduisant le nombre de bits n\'ecessaire au plan secondaire,
  lesquelles se ram\`enent \`a une compression de la configuration
  cellulaire, associ\'ee \`a des m\'ethodes d'acc\`es sp\'ecialis\'ees au codage de
  la compression, renfor\c{c}ant ainsi la tendance des machines
  s\'equentielles \`a l'\'echange de temps contre de l'espace.} %
par cellule pour le calcul d'une g\'en\'eration de \arule{Life}, plus
$c \times  i$ bits pour le codage des instructions, $i$ \'etant la taille
moyenne d'une instruction. %
$c$ et $a$ sont typiquement tous les deux de l'ordre de quelques
centaines d'\bialt{instructions}{cycle} \'el\'ementaire d'une machine
monoprocesseur id\'eale pour une version na\"ive~(non optimis\'ee) de
l'\'emulateur de la machine cellulaire. %
Une impl\'ementation par lookup-table avec collecte pipelin\'ee du
voisinage augmente $c$ et diminue $a$ d'un ordre de grandeur. %

\begin{equation}
  \frac{\text{Space}}{\text{Time}} = \frac{bN+ci+M_{\text{SISD}}}{aN}
  \label{eq-life-sisd}
\end{equation}

Avec:\[b = 2,~i = 32,~c \approx a \approx 10^2,~M_{\text{SISD}}
\approx 10^6\]

La valeur de $a$ et $c$ d\'ependent du nombre $k$ d'\'etats par cellule et
du rayon $r$ du voisinage cellulaire. %
La taille du mot de voisinage~(9~dans le cas de \arule{Life})
d\'etermine ces valeurs. %

L'impl\'ementation de~\arule{Life} sur une machine
SIMD~$M_{\text{SIMD}}$ utilise $a^{\prime}$~cycles de base pour le
calcul d'une g\'en\'eration, ind\'ependamment du nombre de cellules,
et~$b^{\prime}$~bits par cellule, les~$c^{\prime} \times  i^{\prime}$~bits
d'instructions \'etant traditionnellement consid\'er\'es comme n'appartenant
pas stricto sensu \`a la machine SIMD, mais \`a une indispensable machine
s\'equentielle, h\^ote de toute machine SIMD\@. %
Les $a^{\prime}$~cycles de calcul d\'ependent de la machine SIMD choisie,
une machine SIMD typique est plus difficile \`a proposer qu'une
machine SISD id\'eale, sur le GAPP, de l'ordre de la centaine pour une
impl\'ementation a priori de la fonction de transition\footnote{%
  L'algorithme arithm\'etique de \arule{Life}, comptage plus seuil.}, %
de l'ordre du millier pour une impl\'ementation a posteriori de celle-ci\footnote{%
  Une d\'erivation algorithmique du DFA de \arule{Life} calcul\'ee \`a partir
  de la table de transition, laquelle peut \^etre obtenue par observation
  d'une trajectoire sur l'espace des configurations, associ\'ee \`a une
  algorithmique SIMD g\'en\'erique interpr\'etant le DFA sur le mot de voisinage.}. %
De l'ordre d'une dizaine de bits pour les $b^{\prime}$ bits d'espace. %

\begin{equation}
  \frac{\text{Space}}{\text{Time}} = \frac{(b^{\prime}+M_{\text{SIMD}})N}{a^{\prime}}
  \label{eq-life-simd}
\end{equation}

Avec:\[b \approx 10,a^{\prime} \approx 10^2,~\text{SIMD} \approx 10^2\]

L'impl\'ementation de \arule{Life} sur une machine cellulaire pure $M_{\text{CA}}$
utilise $a^{\prime\prime} = 1$ cycle et  $b^{\prime\prime} = 1$ bit par
cellule. %
Par contre le co\^ut li\'e \`a l'architecture s'exprime simplement comme la
taille de la table de transition. %

\begin{equation}
  \frac{\text{Space}}{\text{Time}} =
  \frac{(b^{\prime\prime}+M_{\text{CA}})N}{a^{\prime\prime}}
  \label{eq-life-cell}
\end{equation}

Avec:\[a^{\prime\prime} = 1,~b^{\prime\prime} = 1,~\text{CA} = 512\]

Les valeurs d'espace propos\'ees pour le terme $M_{\text{SISD}}$ et le
facteur $M_{\text{SIMD}}$ dans les \'equations~\ref{eq-life-sisd}
et~\ref{eq-life-simd} repr\'esentent un nombre de transistors CMOS id\'eal
pour ces architectures. %
Le facteur $M_{\text{CA}}$ de l'\'equation~\ref{eq-life-cell} n'est pas
un nombre de transistors approxim\'e, mais un nombre exact de bits. %
Cette diff\'erence de type entre les objets compt\'es se trouve
consid\'erablement pond\'er\'ee par les croissances respectives des tailles
d'espaces n\'ecessaires \`a la r\'ealisation d'une fonction de transition
lorsque la fonction de transition devient un param\`etre des \'equations
pr\'ec\'edentes. %
Sur une machine SISD l'augmentation de la taille du mot de voisinage
induit une croissance lin\'eaire de la taille de l'espace et du temps
utilis\'e pour calculer une trajectoire cellulaire. %
L'augmentation de la taille d'espace ne d\'epend que de la taille du mot
de voisinage, pas de la complexit\'e de la fonction de
transition\footnote{%
  A condition, bien s\^ur, de ne pas transformer la fonction de
  transition en lookup-table.}, %
mais l'augmentation du temps de calcul d\'epend de la taille de la
configuration initiale. %
Sur une machine SIMD l'augmentation de la taille du mot de voisinage
induit \'egalement une croissance lin\'eaire de la taille de l'espace et
du temps utilis\'e pour calculer une trajectoire cellulaire. %
Par contre l'augmentation de la taille d'espace d\'epend de la taille
du mot de voisinage et de la complexit\'e de la fonction de transition,
ce facteur de croissance de la taille d'espace est donc sup\'erieur \`a
celui d'une machine SISD\@. %
De plus la taille de l'architecture est fonction de la taille de la
configuration initiale, ce facteur pr\'esentant \'evidemment la source
principale des besoins de ressource spatiale. %
Par contre, pour prix de cette augmentation lin\'eaire des besoins
d'espace par rapport \`a la version SISD, le temps de calcul n'est plus
fonction de la taille de l'espace initial. %
Sur une machine cellulaire pure l'augmentation de la taille du mot de
voisinage induit une croissance exponentielle de la taille d'espace
via une croissance exponentielle de la taille de la table elle-m\^eme
facteur de la taille de la configuration initiale, le temps se
r\'eduisant \`a l'unit\'e. %

\section{Plan de lecture}

Dans le chapitre~\mygetref{space-time-law} nous pr\'esentons la
structure des objets de notre \'etude:~l'espace, le temps et les lois, %
ainsi que les m\'ethodes applicables sur ces objets:~\'edition, r\'eduction,
transformation, visualisation pour l'espace et le temps, plus
construction, alt\'eration, composition, m\'elange et s\'eparation pour les
lois. %
Nous pr\'esentons \'egalement dans ce chapitre l'architecture de nos
moteurs cellulaires. %

Dans le chapitre~\mygetref{rule-space} nous pr\'esentons les outils
conceptuels utilis\'es pour la quantification de l'espace des r\`egles, et
nous proposons un nouveau param\`etre de contr\^ole pour l'exploration et
la structuration de cet espace:~le param\`etre d'ench\^assement. %

Le chapitre~\mygetref{about-anneal} est consacr\'e \`a une \'etude
exp\'erimentale de la r\`egle~\arule{Anneal} bas\'ee sur l'exploitation de
nos outils. %
Les exp\'eriences r\'ealis\'ees nous am\`enent \`a proposer l'existence d'une
relative invariance spatio-temporelle des trajec\-toires associ\'ees \`a
une transition de phase dans l'espace des configurations initiales. %

La conclusion introduit les limites de notre approche qui \`a associ\'e
un constructivisme bottom-up \`a une \aquote{gestion de tables
  g\'en\'eralis\'ee}. %
Nous \'evoquons pour finir une approche alternative de r\'ealisation
cellulaire purement algorithmique. %

\mydumpfull{anneal-sum-1k-8k}%
{Une superposition de r\'eduction int\'egrale}%
{Une superposition de r\'eduction int\'egrale de \arule{Anneal}}%
{Cette image est obtenue par la superposition de la r\'eduction int\'egrale
  $\sum_{i=0}^{i=n}t_{i}xy$ des~$2^{10}$ et~$2^{13}$ premi\`eres
  g\'en\'erations d'un espace initial de taille~\asquare{1024} et de
  densit\'e\,\Sundemi.}

\mydumpfull{anneal-sum-1k-8k-xv}%
{}%
{Une autre colormap pour la figure
  \protect\mygetdump{anneal-sum-1k-8k}}%
{Une colormap en bandes. %
  Pour cette image, comme pour toutes les images en niveaux de gris
  que nous avons pr\'esent\'ees, un p\'eriph\'erique de sortie capable de
  restitution nuanc\'ee nous fait d\'efaut.}

%%
%% Local Variables:
%% mode: auto-fill
%% iso-accents-minor-mode: nil
%% TeX-command-default: "mytex"
%% TeX-master: "top"
%% Local IspellParsing: "latex-mode ~latin1"
%% Local IspellPersDict: "~/.ispell_lisp"
%% LocalWords: "l'auto-reproduction s\'equencement CA LAB cellsim SunView XView"
%% LocalWords: "Felicity scamper pattern cellang d'AC SIMD wafer scale l'IA MIMD"
%% LocalWords: "matching holographique multim\'edia multi-calculateur ad hoc WAN"
%% LocalWords: "Wide Area Network t\'el\'ecommunication bottom-up von Neumann cpu PE"
%% LocalWords: "adressable repartition-du-calcul l'overhead uds NEW DSP bitblt"
%% LocalWords: "morphog\'en\'etique g\'en\'ericit\'e cardinalit\'e fractales g\'en\'eticit\'e uvre"
%% LocalWords: "it\'er\'ee ergodique analytique-synth\'etique param\'etr\'e Neuman Turing"
%% LocalWords: "auto-reproducteur impl\'ement\'ee Conway univer GAPP it\`ere it\'erant"
%% LocalWords: "lookup-table mapping ind\'eterminisme indiff\'erentiable Langton DFA"
%% LocalWords: "\'epig\'en\'etique morphog\'en\`eses espace-temps Jaqueline Signorini past"
%% LocalWords: "auto-reproductrice monoprocesseur l'\'emulateur impl\'ementation"
%% LocalWords: "stricto sensu SISD CMOS approxim\'e space-time-law rule-space in"
%% LocalWords: "about-anneal Anneal constructivisme"
%% End:

%% Pourquoi notre nouvel environnement?
%% Optimisation d'usage de ressources
%% r\'eflexion morphog\'en\'etique
%% espace, temps et calcul
%% Plan de lecture

%%
%% Local Variables:
%% mode: auto-fill
%% iso-accents-minor-mode: nil
%% TeX-command-default: "mytex"
%% TeX-master: "top"
%% Local IspellParsing: "latex-mode ~latin1"
%% Local IspellPersDict: "~/.ispell_lisp"
%% LocalWords: "impl\'ementatoire SIMD Lattice Gas Automata cytosquelette"
%% End:
